\section{Four forms of loss in lattice gauge reasoning}
\label{sec:failures}

The calculus earns methodological content only if it differentiates failures
that ordinary language treats alike.  We analyze four exact logical forms.
Each begins from a natural estimate and ends with a different repair.

\subsection{Loss of the admissible quantifier domain}

Let $p_t$ be a heat kernel on a compact connected group.  As $t\downarrow0$,
$p_t(g^{-1}\,\cdot)$ concentrates near $g$.  If $g\neq h$, choose disjoint
neighborhoods of the two points.  Concentration then implies
\begin{equation}
\label{eq:tv-translate}
 \left\|p_t(g^{-1}\,\cdot)-p_t(h^{-1}\,\cdot)\right\|_{\mathrm{TV}}
 \longrightarrow 1.
\end{equation}
A worst-translate influence coefficient therefore becomes maximally bad in
the small-time regime where the local measure is most concentrated.

The estimate is correct; its proposed interpretation is not.  Arbitrary
group translates need not correspond to neighboring conditionings inside the
same physically coherent small-field collar.  The representation has exposed
translation and concentration while suppressing the admissibility relation
among boundary states.  The repair is to type the supremum: compare only
conditionings connected by a certified good collar, and treat the complement
through a charged defect or polymer channel.  Replacing the total-variation
constant by a more optimistic local calculation would not repair the missing
quantifier domain.

\subsection{Loss of boundary charge}

Consider one $SU(2)$ link with six incident plaquette staples.  A center-valued
boundary can arrange five identity staples and one $-I$ staple.  For a
heat-kernel plaquette weight, the exact conditional density is proportional
to
\begin{equation}
\label{eq:frustrated-density}
 p_t(U)^5p_t(-U)\,dU.
\end{equation}
Write $\theta\in[0,\pi]$ for the class angle of $U$.  The small-time rate is,
up to the common heat-kernel normalization,
\[
 I(\theta)=5\theta^2+(\pi-\theta)^2.
\]
It is minimized at $\theta=\pi/6$.  The plaquette facing the negative staple
then has angle $5\pi/6$: the one-link conditional is not concentrated on a
configuration in which every adjacent plaquette is small.

The boundary pattern itself carries an action cost.  Conditioning on it and
then demanding a boundary-uniform one-link estimate erases that cost from the
probability accounting.  The repair is a retained boundary charge, a
multi-link integration, or a contour variable that assigns the frustration
to the correct state.  The failed statement was uniform over a state space
too small to remember why the hostile conditioning was rare.

\subsection{Loss of adaptive history}

Fixed intersection moments need not control sequential conditionals.  Let
$0<q<1$ and let Bernoulli variables satisfy
\[
 \mathbb P(X_1=X_2=1)=q^2,
 \qquad
 \mathbb P(X_1=X_2=0)=1-q^2.
\]
Then $\mathbb P(X_j=1)=q^2\leq q$ and
$\mathbb P(X_1=X_2=1)=q^2$, so the one- and two-site intersection bounds
match the familiar powers of $q$.  Nevertheless
\[
 \mathbb P(X_2=1\mid X_1=1)=1.
\]
The intersection data have forgotten which history was revealed before the
next choice.  Adaptive domination needs a filtration-indexed state or a
uniform conditional theorem.  More moments of the same untyped kind do not
automatically recover the history.

\subsection{Loss of normalization}

For a defect event $B$ under boundary state $y$, the exact probability is
\begin{equation}
\label{eq:partition-ratio}
 \mathbb P_y(B)=\frac{Z_B(y)}{Z(y)}.
\end{equation}
An action lower bound on configurations in $B$ controls a numerator.  It does
not control the ratio unless the denominator, or a suitable reference
free-energy difference, is bounded on the same state domain.  A hostile
boundary can suppress $Z(y)$ at the same exponential scale as $Z_B(y)$.
The representation by unnormalized action charge has forgotten entropy and
normalization.  Recovery requires a conditional partition-ratio theorem, not
only a stronger energetic inequality.

\subsection{Why the distinctions matter}

The four repairs enlarge different parts of the representation package:
admissible boundary relations, retained center/topological state, adaptive
filtration, and normalized free energy.  They cannot be ranked on a single
fidelity scale.  This is the practical content of certificate profiles.  A
failed proof becomes a typed diagnosis of which state or theorem is missing,
rather than a generic verdict that the approximation was too crude.
